% Clase 2 --- Distributividad de la unión respecto a la intersección (§3.2).
% Ancla: "y lo único que me va a quedar es esta parte acá dibujo A unión B
% intersectado con C [...] idéntico a hacer primeramente la intersección entre
% A y C [...] Se llama un razonamiento con diagrama de Venn."
% Dos paneles con los mismos tres círculos A, B, C: la región sombreada final
% es idéntica en ambos, lo que muestra de un vistazo que
% (A cup B) cap C = (A cap C) cup (B cap C).
\begin{center}
\begin{tikzpicture}
  % ---------------- Panel (a): (A cup B) cap C ----------------
  \begin{scope}
    \begin{scope}
      \clip (0,-0.75) circle (1.5);
      \fill[figamber!45] (-0.9,0.5) circle (1.5);
      \fill[figamber!45] (0.9,0.5) circle (1.5);
    \end{scope}
    \draw[figline] (-0.9,0.5) circle (1.5);
    \draw[figred, line width=0.9pt] (0.9,0.5) circle (1.5);
    \draw[figsetC] (0,-0.75) circle (1.5);
    \node[figlbl, figblue] at (-2.15,1.55) {$A$};
    \node[figlbl, figred] at (2.15,1.55) {$B$};
    \node[figlbl, color=figamber, fill=white, inner sep=1.5pt] at (0,-2.35) {$C$};
    \node[figlbl] at (0,-2.95) {$(A\cup B)\cap C$};
    \node[figlblaux] at (0,-3.45) {(a)};
  \end{scope}

  % ---------------- Panel (b): (A cap C) cup (B cap C) ----------------
  \begin{scope}[xshift=6.2cm]
    \begin{scope}
      \clip (0,-0.75) circle (1.5);
      \fill[figamber!45] (-0.9,0.5) circle (1.5);
      \fill[figamber!45] (0.9,0.5) circle (1.5);
    \end{scope}
    \draw[figline] (-0.9,0.5) circle (1.5);
    \draw[figred, line width=0.9pt] (0.9,0.5) circle (1.5);
    \draw[figsetC] (0,-0.75) circle (1.5);
    \node[figlbl, figblue] at (-2.15,1.55) {$A$};
    \node[figlbl, figred] at (2.15,1.55) {$B$};
    \node[figlbl, color=figamber, fill=white, inner sep=1.5pt] at (0,-2.35) {$C$};
    \node[figlbl] at (0,-2.95) {$(A\cap C)\cup(B\cap C)$};
    \node[figlblaux] at (0,-3.45) {(b)};
  \end{scope}
\end{tikzpicture}\\[0.5em]
{\small\itshape Mismos tres conjuntos $A$, $B$, $C$ en los dos paneles, y
\textbf{la misma región queda sombreada}: (a) se arma tomando primero
$A\cup B$ y recortándola con $C$; (b) se arma tomando por separado $A\cap C$ y
$B\cap C$ y uniéndolas. Que el resultado visual coincida es exactamente lo que
prueba la distributividad, sin necesitar la demostración formal por doble
inclusión.}
\end{center}
